% Passes 5308--5313: D4 triality, tomotope quotient, tesseract near-miss, and Latin/Q4 refinement.
\subsection{The Hoffman $576$ stabilizer, $D_4$ triality, and the tomotope quotient}
\label{sec:hoffman-d4-triality-tomotope}

The order-$576$ setwise stabilizer $H$ of the $13$-cell Hoffman cover has a
canonical characteristic subgroup
\[
 N=\langle h\in H:\ |h|\text{ is a power of }2\rangle,
 \qquad |N|=192.
\]
Machine reconstruction gives
\[
 N\cong 2^3\!:\!S_4\cong W(D_4),
 \qquad Z(N)\cong C_2,
 \qquad H/N\cong C_3.
\]
Moreover $N$ contains exactly three normal elementary-abelian $2^3$ subgroups,
and an element representing the quotient $C_3$ cyclically permutes these three
subgroups.  Thus the factor $576/192=3$ is an explicit $D_4$ triality action,
not merely a numerical factorization.

The action of $N$ on the $13$ Hoffman cover cells has orbit structure
\[
 1+12,
\]
and its kernel is exactly $Z(N)$.  On the $12$ moving cells the induced
order-$96$ permutation group is, after the explicit relabeling
\[
 (0,1,\ldots,11)\mapsto(0,3,9,5,6,10,1,4,11,2,8,7),
\]
identical to the group generated by the four published tomotope permutations.
Consequently
\[
 \boxed{W(D_4)/Z(W(D_4))\cong\Gamma(T)}
\]
not only abstractly, but as the displayed degree-$12$ permutation action.
Here
\[
 |\Gamma(T)|=96,
 \qquad \Gamma(T)\cong 2^4\!:\!S_3.
\]
The thirteenth fixed Hoffman cell is not identified with a tomotope face by this
argument.

\paragraph{The tesseract order-$192$ near miss.}
The full signed-permutation symmetry group of the $4$-cube has order $384$ and
contains two natural index-two subgroups of order $192$.  The determinant-$+1$
subgroup $R$ is the orientation-preserving tesseract group, while the even-sign
subgroup is $W(D_4)$.  They are not isomorphic: their element-order censuses are
\[
 R:\quad 1^1\,2^{43}\,3^{32}\,4^{36}\,6^{32}\,8^{48},
\]
\[
 W(D_4):\quad 1^1\,2^{43}\,3^{32}\,4^{84}\,6^{32}.
\]
In particular, $R$ has $48$ elements of order $8$ and $W(D_4)$ has none.
Nevertheless both centers have order two and both central quotients are
isomorphic to
\[
 2^4\!:\!S_3\cong\Gamma(T).
\]
Thus the tomotope is a common order-$96$ central quotient of two distinct
order-$192$ double covers.  The Hoffman stabilizer selects the $W(D_4)$ cover,
not the tesseract rotation cover.

\paragraph{Latin and toroidal-$Q_4$ refinement.}
For the fixed $4\times4$ toroidal-knight graph used in this project, which is
isomorphic to $Q_4$, enumerate all $576$ Latin squares of order four and count
the knight edges whose endpoints carry the same symbol.  The exact census is
\[
 \boxed{576=96+192+288},
\]
where the three classes have respectively $0$, $8$, and $16$ same-symbol
knight edges.  Under the common board symmetry preserving both the row/column
grid and this fixed $Q_4$ edge set, together with independent symbol relabeling,
the $576$ squares split further into orbit sizes
\[
 48,\ 48,\ 96,\ 192,\ 192.
\]
This is an exact object-level statistic on the shared $16$-cell substrate.  The
coincidence of the class sizes $96,192,288$ with the tomotope, $D_4$/tesseract,
and Hoffman-central-quotient scales is recorded as a structural target, not as
a group-isomorphism theorem.

Finally, the Pass~5300 affine-module isomorphism
\[
 H/Z(H)\cong L_+
\]
with the even Klein-Latin parastrophe group $L_+$ sends the distinguished normal
index-three subgroup $W(D_4)/Z$ to the Klein autotopy subgroup of order $96$.
The corresponding quotient $C_3$ is therefore $D_4$ triality on the Hoffman
side and cyclic row/column/symbol parastrophe on the Latin side.  This is an
affine-module statement: the natural degree-$12$ Hoffman and Latin label actions
have nonisomorphic point stabilizers and are not permutation-equivalent.
